\section{Protocol Robustness and Integrity}
\label{sec:robust}

\subsection{The parser assumes the line is hostile}
Over-length input is bounded by the 64-byte frame limit and answered
\texttt{NACK:TOOLONG}, with the remainder of the offending line discarded so
that its tail is not then parsed as a fresh command. Non-printable bytes are
dropped as they arrive rather than buffered. A partial frame is abandoned after
2000\,ms, so a client that dies mid-line releases the parser without
intervention. A bare newline is ignored rather than answered.

An unrecognised command is \texttt{NACK:BADCMD}. The index on a \texttt{MOVE}
is exactly one digit, so \texttt{MOVE:9} is a well-formed request for a cell
that does not exist and is answered \texttt{NACK:RANGE}, while
\texttt{MOVE:99}, \texttt{MOVE:-1} and \texttt{MOVE:} are malformed and are
answered \texttt{NACK:BADCMD}. The two are distinguished because they indicate
different client faults, and both were confirmed against the board.

Every rejection is explicit. The server never answers silence, so a client can
always distinguish a refused request from a lost one by timeout alone.

\subsection{The checksum detects corruption, not forgery}
An optional \texttt{*HH} suffix carries an XOR of every byte before the
\texttt{*}, as two hex digits in either case. When the suffix is absent the
message is accepted unchecked; when it is present and wrong the message is
answered \texttt{NACK:BADCRC} and dropped.

It detects single-byte corruption and most burst corruption on the wire, which
is what a short UART link is likely to suffer. It is not a cryptographic check:
an XOR is trivially forgeable, it cannot detect a transposition of two bytes,
and it offers nothing against replay or against a client that is simply lying
about whose turn it is. Integrity against a hostile client is provided by server
authority, as argued in Section~\ref{sec:authority}, and not by the checksum.
